\section{Local inversion with nonnegligible remote edge mass}
\label{sec:localized-inversion}
Theorem~\ref{thm:LLT} gives a sufficient global smallness condition on
extracted edges. A central local limit needs less: the extracted measure
must have negligible density and negligible low-frequency convolution
on the central window. Its total variation need not tend to zero.
The following refinement retains the raw datum and the density norm.

Fix an even real function $\chi\in C_c^\infty(\R^4)$ equal to one on
the unit ball and supported in the ball of radius two. Let $a_n\to\infty$
with $a_n=o(n^{1/6})$, put $B_n=a_n/\sqrt n$, and set
$\chi_n(\omega)=\chi(\omega/B_n)$. For all sufficiently large $n$ its
support fits in the fundamental torus interval in the first three
coordinates. With the transform convention already fixed, its inverse
kernel on $\Z^3\times\R$ is
\begin{equation}\label{eq:localized-kernel}
 K_n(k,t)=B_n^4\check\chi(B_nk,B_nt),\qquad
 |K_n(k,t)|\le C_M B_n^4(1+B_n|(k,t)|)^{-M}.
\end{equation}
The check denotes the inverse transform on $\R^4$, including its
$(2\pi)^{-4}$ factor. Convolution here and below uses counting measure
in the discrete coordinates and Lebesgue measure in the roof coordinate.
For fixed $K$, let $W_{n,R,K}=\{z:|z-nm_R|\le K\sqrt n\}$.

\begin{theorem}[A local edge criterion]\label{thm:local-edge-LLT}
Assume the small-frequency hypothesis (i) of Theorem~\ref{thm:LLT}.
Let $E_{n,R}$ be finite signed measures with densities $e_{n,R}$ and
suppose $H_{n,R}=\Phi_{n,R}-\widehat E_{n,R}$ belongs to
$L^1(\T^3\times\R)$. Assume, uniformly in $R$, that
\begin{equation}\label{eq:local-tail}
 \int |1-\chi_n(\omega)|\,|H_{n,R}(\omega)|\dd\omega=o(n^{-2}).
\end{equation}
For every fixed central window suppose also that
\begin{equation}\label{eq:local-edge-errors}
 n^2\mathop{\rm ess\,sup}_{W_{n,R,K}}|e_{n,R}|\longrightarrow0,
 \qquad
 n^2\sup_{W_{n,R,K}}|K_n*E_{n,R}|\longrightarrow0.
\end{equation}
Then the raw density conclusion \eqref{eq:LLT} holds on that window.
There is no small-total-variation hypothesis on $E_{n,R}$.
\end{theorem}
\begin{proof}
As in the proof of Theorem~\ref{thm:LLT}, integrability of $H$ first
identifies a continuous density for $\eta-E$, and adding $e$ gives
an actual mixed density for $\eta$. Splitting this identity with
$\chi_n$ gives the exact equality almost everywhere
\begin{equation}\label{eq:exact-local-decomposition}
 p_{n,R}=K_n*\eta_{n,R}
       +(e_{n,R}-K_n*E_{n,R})
       +\mathcal F^{-1}[(1-\chi_n)H_{n,R}].
\end{equation}
All terms are defined: $E$ is finite, $K_n$ is bounded and integrable,
and the last inverse transform is absolutely convergent.
The last term is at most $(2\pi)^{-4}$ times the integral in
\eqref{eq:local-tail}. The middle term is controlled by
\eqref{eq:local-edge-errors}.

For the first term, its Fourier integral has the factor $\chi_n$.
Insert the assumed small-frequency expansion and use
$v=\sqrt n\,\omega$. The support has $|v|\le2a_n$ and
$nO(|\omega|^3)=O(a_n^3/\sqrt n)=o(1)$. The uniform quadratic
bound on $\lambda_R$ gives a Gaussian dominator, while the amplitude
converges to one. The exponentially small spectral remainder, after
integration and multiplication by $n^2$, is bounded by
$Ca_n^4\rho^n$ and tends to zero. Truncate first at a fixed $|v|$,
then let that truncation grow. This proves uniformly in $R$ and the
central variable that $n^2K_n*\eta_{n,R}$ converges to the Gaussian
in \eqref{eq:LLT}. Combining the three terms proves the assertion.

The kernel is only an analytical device in the exact decomposition
\eqref{eq:exact-local-decomposition}. The two correction terms are
retained and estimated; the result is not a local limit for a
smoothed replacement of the observed law.
\end{proof}

\begin{corollary}[Remote mass need not be small]\label{cor:remote-edges}
Suppose $\|E_{n,R}\|_{\TV}\le C$ and, for the window under
consideration, the support of $E_{n,R}$ is at distance at least
$\varepsilon\sqrt n$ from $W_{n,R,K}$, for a fixed
$\varepsilon>0$. Then \eqref{eq:local-edge-errors} holds.
In particular an extracted density of bounded, nonvanishing mass
can satisfy the local edge hypotheses.
\end{corollary}
\begin{proof}
The density vanishes almost everywhere in the central window. For
$z$ in that window, \eqref{eq:localized-kernel} gives
\[
 n^2|(K_n*E_{n,R})(z)|
 \le CC_M a_n^4(1+\varepsilon a_n)^{-M}\longrightarrow0
\]
by taking $M>4$. The Schwartz bound is available for every fixed $M$
because $\chi$ is smooth. These constants are kernel constants, not
the derivative-growth constants of the billiard in Lemma~\ref{lem:splice}.
\end{proof}

The same conclusion holds for a sum of remote extracted terms whose
total variations are summable with bounded total, and a remaining
nearby extracted density satisfying \eqref{eq:local-edge-errors}.
This follows by linearity and the triangle inequality. Thus individual
edge coefficients do not all have to be $o(n^{-2})$ merely because
one proves a central limit. Nevertheless, one must still construct
the full extraction and prove \eqref{eq:local-tail}. Neither this
corollary nor the single-word estimate in Proposition~\ref{prop:general-edge}
controls the sum of all central critical words or singular boundaries.
A restriction of a density to a sharp window can itself create new
jump terms; it is not automatically an integrable residual.

For weighted measures the same proof gives the corresponding amplitude
times the Gaussian, provided the small-frequency expansion, residual
estimate and local edge errors are verified for that weighted class.
This is the precise additional input needed before applying
Proposition~\ref{prop:conditioning} to a proposed family of insertions.
